\section*{Patristic and Saintly Reception in Concert}

The witnesses do not supply detached ornaments to the Sunday. Read across the
formulary, they disclose a coherent ascent from covenantal poverty to
sacramental communion and from mercy received to a life made Eucharistic.

\begin{longtable}{>{\raggedright\arraybackslash}p{0.17\textwidth}>{\raggedright\arraybackslash}p{0.35\textwidth}>{\raggedright\arraybackslash}p{0.40\textwidth}}
\toprule
\textbf{Sense} & \textbf{Principal witnesses} & \textbf{Contribution to the whole formulary}\\
\midrule
\endhead
Literal and historical & Lyra, Theodoret, Cassiodorus, Bellarmine, Jerome,
Tertullian, Cyril, Bede & Israel pleads covenant amid sanctuary desolation;
the good Law serves promise without giving life; Jesus honors priestly
examination and freely cleanses all ten; manna remains the Father's historical
gift to Israel before it becomes Christian sacramental typology.\\
\addlinespace
Christological and ecclesial & Eusebius, Augustine, Cassiodorus, Tertullian,
Rabanus, Albert, Aquinas, Bonaventure & Christ is Abraham's Seed, covenantal
inheritance, true Priest and temple, eternal refuge, crucified speaker of the
entrusted times, and heavenly bread. Those incorporated into him become the
one seed, a reconciled people gathered into Eucharistic thanksgiving.\\
\addlinespace
Moral and penitential & Origen, Augustine, Gregory the Great, Bede,
St~Anthony, Bruno, Albert, Bonaventure, Newman & Poverty renounces self-made
justice; the Law reveals the wound; Christ commands a faith that walks before
sight; confession restores ecclesial communion; gratitude attributes every
good to God; providential trust offers every season rather than submitting to
chance or fear.\\
\addlinespace
Anagogical & Augustine, Jerome, Cassiodorus, Albert, Aquinas, Bonaventure,
St~John Paul~II & Ruined sanctuary opens toward heavenly inheritance; passing
generations seek refuge in the eternal Word; viaticum pledges future fruition;
the repeated \latin{augméntum} names the pilgrim's growth until redemption is
possessed without loss.\\
\bottomrule
\end{longtable}

\subsection*{The lepers, confession, and the Church's ministry}

The penitential reading is not an isolated medieval ingenuity. Augustine's
\emph{Quaestiones evangeliorum} already joins the priestly showing to the
Church's ministry of doctrine and sacrament. Gregory the Great uses the
mixture of healthy and diseased skin for the doctrinal mixture of truth and
error. Bede inherits this ecclesial horizon while insisting that faith's
living operation receives the saving word. The anonymous \emph{De vera et
falsa poenitentia}, Gratian, and Peter Lombard receive the sequence as a
major Western locus for contrition and confession.

The doctors refine rather than dissolve the type. Bruno of Segni distinguishes
God's remission from the priest's judgment. St~Anthony describes contrition,
confession, and satisfaction under \latin{ite---osténdite---dum irent}.
Albert states \latin{Deus enim solus emundat peccata}; Bonaventure gives
interior intention its true place before completed confession. At the same
time Catholic doctrine affirms more than external verification: by Christ's
delegated power, the ordained priest truly absolves and reconciles in a
judicial act. The type therefore holds together divine forgiveness, conversion
of heart and sacramental encounter, judgment and restoration to communion.

This reception also carries an ascetical demand. Christ's question concerning
the nine rebukes culpable ingratitude; healing is gratuitous, but praise is
owed. Augustine directs the rebuke inward: do not attribute amendment to self.
St~Anthony calls the grateful man the keeper of the gift. Albert's ordered
verbs---recognition, return, praise, prostration, humility, thanksgiving---show
how grace matures into worship. The leper type, properly used, names the
worshipper's sin, error, and estrangement and the Lord's restoring mercy. It
cannot be turned outward as a diagnosis of Hansen disease, disability, or
visible difference.

\subsection*{The one Samaritan, the nations, and Israel's promise}

Augustine, Bede, Bruno, St~Anthony, and Bonaventure find spiritual significance
in ten, nine, and one. Ten can represent humanity
measured by the Decalogue and universally enclosed under sin. Nine can signify
incompletion; one the unity in which grace is kept through praise. The
Samaritan figures the nations' firstfruits, a stranger brought near who
confesses the God of Israel at Jesus' feet. Titus of Bostra, preserved through
the \emph{Catena aurea}, reads Christ's passage as peace between estranged
peoples, joining two into one new man.

\Needspace{5\baselineskip}
This traditional Gentile-inclusion type belongs naturally with Galatians: the
promised Seed is Christ, and in him the blessing reaches every nation. It also
belongs with the Alleluia's plural \latin{nobis}, the refuge of successive
generations, and the Communion's grateful \latin{nobis}. Theophylact gives its
moral law: alien birth excludes no one from pleasing God, and sacred ancestry
gives no right to pride.

The type is strongest when its scriptural order is kept. Israel receives the
promise, Law, priesthood, psalms, and manna; Israel gives the Messiah according
to the flesh; through Israel the nations glorify God in Christ. Luke does not
say that all nine are Jews or declare a collective rejection. Thus the
traditional figure of the nations can be preached positively under the
Catholic confession that Israel's gifts and calling are irrevocable and that
the Jews may not be represented as rejected or accursed (\emph{Nostra
aetate}~4).

\subsection*{From spiritual manna to Eucharistic viaticum}

The Wisdom reception also grows by ordered development. Origen first hears the
one Logos proportioning nourishment to the soul. Gregory of Nyssa joins that
condescension to the Virgin-born Word. Gregory the Great sees the one divine
discourse addressing differing hearers. These direct uses establish a
christological and scriptural grammar: heavenly Wisdom is one, yet gives what
leads each receiver toward maturity.

Rabanus then comments directly from Israel's manna through John~6 to Christ's
flesh and blood. Albert, Aquinas, and Bonaventure articulate the Eucharistic
mystery as heavenly, unitive, life-giving, abundant, and viaticum. The
\emph{Roman Catechism} names Communion's increase of grace and charity; St~John
Paul~II names pilgrim food and messianic foretaste. Gregory's scriptural use
complements rather than disciplines this Eucharistic typology: the same divine
condescension that speaks to human capacity gives the one Christ sacramentally
and enlarges the receiver for glory.

\subsection*{Recent papal reception: cry, walk, return, worship}

Benedict~XVI distinguishes bodily health from definitive salvation and places
faith and gratitude at the hinge. Reading Naaman with the Samaritan, he presents
universal salvation as coming through Israel and reaching fulfilment in
Christ. Across several occasions Francis follows the Gospel's movement from
the cry and obedient walk to return and thanksgiving. He also
joins the Samaritan's gratitude to \emph{Eucharistia}, the Church's thanksgiving.
Leo~XIV, preaching the pericope in 2025, emphasizes both the freedom of Christ's
healing and the one man's recognition of salvation, while refusing the use of
``leper'' as a label for outsiders. These modern papal receptions do not
replace the patristic senses; they show their durable pastoral form.

At the end, the line of interpretation remains one: covenant mercy elicits a
cry; grace commands and enables a road; healing seeks recognition; recognition
returns as worship; thanksgiving enters sacrifice and sacrament; heavenly food
sends the reconciled pilgrim onward. The Fathers, doctors, saints, and recent
pastors speak in different idioms, but each serves this movement from gift to
grace-enabled return.
